\section{Limitations}
\label{sec:limitations}

We ship a v0.1 draft. Five named limitations bound the scope of the
claims above. Carried forward verbatim from the project's narrative
report (NARRATIVE\_REPORT.md, §6) as the spine of honesty.

\paragraph{L1: Bag-of-tokens vs.\ semantic embedding ---
\textit{partially addressed in v0.2}.} The v0.1 default embedder is
SHA1-hashed bag-of-tokens and cannot separate framing words from
content words in encoding-task prompts. We added a pluggable embedder
interface and a sentence-transformers
(\texttt{all-MiniLM-L6-v2}) wrapper \cite{reimers2019sbert}; E1b
(\cref{sec:exp:e1b}) demonstrates that the true forward pass with this
real semantic embedder recovers the human SRE Cohen's $d$ within the
Symons \& Johnson 1997 reference band. The bag-of-tokens default
remains for backward compatibility and dependency-free testing; it is
documented as inadequate for forward-pass SRE work.

\paragraph{L2: $\piself$ update rule ---
\textit{partially addressed in v0.4}.}
\Cref{eq:pi} now carries a documented variational interpretation
(\cref{sec:method:self}): the formula is the simplest calibrated
reliability link recovering, in the large-window limit, the
regularised inverse-variance shape of a Gamma--Inverse-Gamma
conjugate update under Gaussian observation noise. The
hyperparameter $\gamma$ is exposed as the explicit
\texttt{prior\_precision} kwarg in code. We do \emph{not} claim
\cref{eq:pi} is the exact closed-form posterior of a fully specified
generative model; a strict derivation that algebraically produces
\cref{eq:pi} from a single likelihood/prior pair is named as v0.5
work. The qualitative claims of the paper (Sections 4.2--4.5) do
not depend on which exact form of $\piself$ is used --- only on the
property that $\piself \to 1$ at low PE variance and $\to 0$ at
high PE variance.

\paragraph{L3: Human-data fit ---
\textit{partially addressed in v0.2}.} The model's predicted Cohen's
$d$ between self-ref and semantic conditions now meets the Symons \&
Johnson 1997 meta-analytic reference band (E1b, \cref{fig:e1b-cohensd}).
However, we have not yet:
(a) fit our model to Rogers, Kuiper \& Kirker \cite{rogers1977sre} raw
per-subject recall data, or
(b) compared head-to-head with CMR3 \cite{cohen2022cmr3} on emotional
memory benchmarks (E5 in \cref{sec:exp:future}),
(c) attempted clinical fits against attenuated SREs in dissociation
and depression (E6).
CMR3 reimplementation is a multi-day effort we have not budgeted.

\paragraph{L4: Deterministic experiments without noise.} All three
experiments are deterministic given seed; cell-level variance is exactly
zero in E1. Adding realistic encoding noise --- Gaussian on forget\_score,
or stochastic embedding perturbation --- is necessary for honest variance
estimates and for the OLS regressions in E2 to be applied to data with
the variance structure they assume. We expect the qualitative
conclusions to be robust but cannot claim it from current evidence.

\paragraph{L5: Retrieval uses encoding-time $\muself^{\text{enc}}$ ---
\textit{addressed in v0.4}.} The retrieval ranking (\cref{eq:rank})
now prefers $\cos(\muself^{\text{enc}}, \muself^{\text{now}})$ over
the looser $\cos(e_i, \muself^{\text{now}})$, where
$\muself^{\text{enc}}$ is the agent's self prior at the moment memory
$i$ was formed --- persisted in a new \texttt{mu\_self\_at\_encoding}
column on the memory store. This makes the model strictly faithful to
Tulving's encoding-specificity principle
\cite{tulving1973encodingspec}: the agent matches the self it has now
against the self it had then, not against the literal content of the
memory. A graceful fallback to the embedding path is retained for
legacy rows that pre-date the snapshot column.

\paragraph{Path forward.}
v0.2 addressed L1 (SBert wrapper + E1b forward pass) and the first
half of L3 (Cohen's $d$ in the S\&J reference band). v0.4 addresses
L2 (variational derivation of $\piself$) and L5 (encoding-time
$\muself^{\text{enc}}$ persistence). The remaining items are L4
(more thorough noise modelling than the $\sigma{=}0.05$ Gaussian we
already add at the forget-score level), the CMR3 head-to-head
reimplementation, and the clinical E6 attenuation fit. None of these
change the central mechanism; they sharpen its quantitative claims
and broaden empirical coverage.
